%% EXEMPLO FICTÍCIO — Atividade 1: eletrônica e placa de circuito impresso
\begin{atividade}{Placa de potência do estimulador do NEBULA, revisão B}{
  tipo         = pcb,
  modalidade   = COMP,
  alternativa  = PES,
  horas        = 150,
  inicio       = 2024-03,
  fim          = 2024-09,
  projeto      = {NEBULA, estimulação elétrica funcional para ciclismo},
  papel        = Responsável,
  supervisor   = {Rafael Fictício, supervisor do projeto NEBULA},
  registros    = {NBL-114; NBL-121; NBL-133; NRO-PRO-014-901; NRO-PRO-003-901},
  competencias = {I, II, III, V},
  disciplinas  = {Eletrônica Analógica; Eletrônica de Potência; Circuitos Elétricos II; Compatibilidade Eletromagnética},
}

\begin{contexto}
O NEBULA é o projeto da equipe que leva uma pessoa com lesão medular a pedalar uma bicicleta
adaptada por meio de estimulação elétrica funcional (FES): pulsos de corrente, aplicados por eletrodos
de superfície, contraem na ordem certa os músculos das pernas. O coração do sistema é o estimulador,
e a placa de potência é a parte dele que gera os pulsos.

A revisão A da placa funcionava, mas os testes de 2023 (\texttt{NRO-PRO-003-900}) registraram quatro
problemas: a amplitude da corrente caía quando a impedância do eletrodo subia; o estágio de saída
chegava a \SI{71}{\celsius} em estimulação contínua; o chaveamento da saída induzia ruído no sinal do
encoder do pedal; e não havia limite de corrente por hardware --- a proteção dependia só do firmware.

Os requisitos vêm da USRS do projeto (\texttt{NRO-PRO-004-901}): pulsos bifásicos com carga
balanceada; amplitude de 0 a \SI{100}{\milli\ampere}, em passos de \SI{1}{\milli\ampere}; largura de
50 a \SI{500}{\micro\second}; frequência de 20 a \SI{50}{\hertz}; quatro canais; e isolação galvânica
entre o usuário e todo ponto ligado à rede elétrica (o carregador e o USB). A revisão B tinha de
resolver os quatro problemas sem sair desses requisitos.
\end{contexto}

\begin{contribuicao}
Fui a responsável pela revisão B, do esquemático à placa montada e testada:
\begin{itemize}
  \item reprojetei o estágio de saída (a fonte de corrente e a ponte H) e criei o limite de corrente
    por hardware, que não existia;
  \item escolhi os componentes críticos --- o amplificador operacional, os MOSFETs de \SI{250}{\volt}, os
    isoladores --- e registrei as escolhas e as alternativas no design record (\texttt{NRO-PRO-014-901});
  \item fiz o layout das quatro camadas, montei as duas primeiras placas no LABBIO e conduzi a primeira
    energização.
\end{itemize}
Não foram meus: o firmware que gera os pulsos (Lucas Fictício); o conversor boost, que veio da
revisão A e em que só ajustei a tensão de saída; e a revisão do layout antes da fabricação, feita pelo
supervisor do projeto.
\end{contribuicao}

\begin{desenvolvimento}
\fichatecnica{
  ferramenta  = {KiCad 8.0},
  revisao     = {B (a revisão C incorpora a errata desta)},
  camadas     = {4 camadas, FR-4 de 1,6\,mm},
  dimensoes   = {100 × 80\,mm},
  componentes = {146},
  alimentacao = {Bateria 3S (11,1\,V); DC-DC isolado de 3\,kV para o lado do usuário; boost de 120\,V},
  fabricacao  = {Fabricação terceirizada; montagem manual no LABBIO (duas placas)},
  normas      = {IEC 60601-1 e IEC 60601-2-10 como referência de projeto, sem certificação},
  arquivos    = {Repositório \texttt{nebula-hw}, tag \texttt{rev-b}},
}

\textbf{A arquitetura.} A placa tem dois lados, separados por uma barreira de isolação
(Figura~\ref{fig:ex-blocos}). Do lado ligado ao USB e ao carregador ficam o microcontrolador e a
bateria; do lado do usuário, o conversor boost que gera os \SI{120}{\volt}, os conversores
digital-analógicos que definem a corrente de cada canal e os estágios de saída. Só atravessam a barreira
a energia, por um conversor DC-DC isolado, e os sinais digitais, por isoladores. A revisão A não tinha
essa separação: o USB e o eletrodo dividiam o mesmo terra.

\begin{figure}[htbp]
  \centering
  \resizebox{\linewidth}{!}{% Diagrama de blocos da placa de potência, revisão B (exemplo fictício)
\begin{tikzpicture}[
    bloco/.style={draw=nro-tinta, line width=0.5pt, fill=white, rounded corners=1.2pt,
      minimum height=11mm, text width=29mm, align=center, font=\footnotesize, inner sep=1.5pt},
    critico/.style={bloco, fill=nro-synapse!18},
    seta/.style={-{Stealth[length=1.8mm]}, line width=0.55pt, draw=nro-tinta},
    alta/.style={seta, draw=nro-pulso-texto, line width=0.8pt},
    x=1mm, y=1mm]
  % colunas: controle | barreira | lado do usuário
  \node[bloco]   (mcu)   at (0,32)  {Microcontrolador\\{\scriptsize firmware, USB}};
  \node[bloco]   (bat)   at (0,0)   {Bateria 3S\\{\scriptsize 11,1\,V, 2,6\,Ah}};
  \node[bloco]   (iso)   at (44,32) {Isoladores digitais\\{\scriptsize SPI e habilitações}};
  \node[bloco]   (dcdc)  at (44,0)  {DC-DC isolado\\{\scriptsize 3\,kV, 12\,V}};
  \node[bloco]   (dac)   at (88,32) {DAC de 12 bits\\{\scriptsize a referência de corrente}};
  \node[bloco]   (boost) at (88,0)  {Conversor boost\\{\scriptsize 12\,V para 120\,V}};
  \node[critico] (fonte) at (132,32) {Fonte de corrente\\{\scriptsize amp. op. e MOSFET}};
  \node[critico] (ponte) at (132,0)  {Ponte H\\{\scriptsize pulso bifásico}};
  \node[critico] (lim)   at (88,16) {\scriptsize Limite de corrente: 110\,mA, trava e parada de emergência};
  \node[bloco, text width=22mm] (ele) at (132,-17) {Eletrodos\\{\scriptsize 4 canais}};
  \draw[seta] (mcu) -- (iso);
  \draw[seta] (iso) -- (dac);
  \draw[seta] (dac) -- (fonte);
  \draw[seta] (bat) -- (dcdc);
  \draw[seta] (dcdc) -- (boost);
  \draw[alta] (boost) -- node[above, font=\scriptsize, text=nro-pulso-texto] {120\,V} (ponte);
  \draw[seta] (fonte) -- (ponte);
  \draw[seta] (ponte) -- (ele);
  \draw[seta] ([xshift=-10mm]fonte.south) |- ([yshift=2mm]lim.east);
  \draw[seta] ([yshift=-2.5mm]lim.east) -| ([xshift=-7mm]ponte.north);
  \draw[seta] (bat.north) -- (mcu.south);
  % a barreira de isolação
  \draw[dashed, line width=0.7pt, nro-pulso-texto] (44,-9) -- (44,43)
    node[above, font=\scriptsize, text=nro-pulso-texto] {barreira de isolação (3\,kV)};
  \node[font=\scriptsize, text=nro-cinza, anchor=west] at (-14,41) {lado ligado ao USB e ao carregador};
  \node[font=\scriptsize, text=nro-cinza, anchor=east] at (146,41) {lado do usuário};
\end{tikzpicture}
}
  \caption{Diagrama de blocos da placa de potência, revisão B. Em verde-claro, os blocos refeitos nesta
    revisão.}
  \fonte{Elaborado pela autora.}
  \label{fig:ex-blocos}
\end{figure}

\textbf{A fonte de corrente.} Cada canal tem uma fonte de corrente controlada por tensão
(Figura~\ref{fig:ex-fonte}): o amplificador operacional força sobre o resistor de medição $R_S$ a
tensão que o DAC define, e a mesma corrente atravessa a carga. Com ganho de malha alto,
\begin{equation}
  I_\text{saída} = \frac{V_\mathrm{DAC}}{R_S}.
  \label{eq:ex-fonte}
\end{equation}
Com $R_S = \SI{10}{\ohm}$, a faixa de 0 a \SI{100}{\milli\ampere} pede $V_\mathrm{DAC}$ de 0 a
\SI{1}{\volt}; o DAC de 12 bits dá passos de \SI{24}{\micro\ampere}, bem abaixo do passo de
\SI{1}{\milli\ampere} do requisito. Escolhi um resistor de 0,1\,\% e 25\,ppm/°C: a exatidão da
corrente é a dele.

A queda de amplitude da revisão A tinha duas causas, que a simulação da Atividade~2 separou: a
tensão de \SI{90}{\volt} não bastava para manter \SI{100}{\milli\ampere} em eletrodos de impedância
alta, e o laço oscilava na borda do pulso. A revisão B sobe a tensão para \SI{120}{\volt} e tem o
capacitor de compensação $C_c = \SI{22}{\pico\farad}$, que a simulação dimensionou.

\begin{figure}[htbp]
  \centering
  % A fonte de corrente de um canal (exemplo fictício): o amplificador
% operacional força sobre R_S a tensão do DAC; a corrente na carga é V_DAC/R_S.
\begin{circuitikz}[font=\footnotesize, line width=0.5pt]
  \draw (0,0) node[op amp, yscale=-1] (oa) {};
  \node at ([xshift=-1.5mm]oa.center) {\scriptsize U1};
  % a referência na entrada não inversora (em cima)
  \draw (oa.+) to[short, -o] ++(-1.4,0) node[left] {$V_\mathrm{DAC}$ (0 a 1\,V)};
  % a saída, o resistor de gate e o MOSFET
  \draw (oa.out) to[R, l=$R_G$, -*] ++(2.4,0) coordinate (g);
  \draw (g) node[nigfete, anchor=G] (q) {};
  \node[right, font=\scriptsize] at (q.east) {Q1 (250\,V)};
  \draw (q.D) -- ++(0,0.5) to[R, l_=$Z_\mathrm{carga}$] ++(0,2.1) node[above] {$V_\mathrm{HV} = 120$\,V};
  % o resistor de medição e a realimentação, por baixo
  \draw (q.S) -- ++(0,-1.9) coordinate (s) to[R, l_={$R_S = 10\,\Omega$}, *-] ++(0,-2) node[ground] {};
  \coordinate (r0) at ([xshift=-9mm]oa.-);
  \draw (s) -- (s -| r0) |- (oa.-);
  \draw (s) -- ++(1.4,0) node[right, align=left, font=\scriptsize] {ao comparador\\do limite (110\,mA)};
  % a compensação, entre a saída e a entrada inversora
  \draw ([xshift=-4mm]oa.-) node[circ] {} -- ++(0,-0.9) coordinate (c1);
  \draw (c1) to[C, l_=$C_c$] (c1 -| oa.out) -- (oa.out);
\end{circuitikz}

  \caption{A fonte de corrente de um canal. A tensão em $R_S$ também alimenta o comparador do limite de
    corrente.}
  \fonte{Elaborado pela autora.}
  \label{fig:ex-fonte}
\end{figure}

\textbf{A dissipação.} No pior caso --- carga em curto ---, o MOSFET dissipa toda a tensão:
\begin{equation}
  P_\mathrm{pico} = V_\mathrm{HV}\, I = \SI{120}{\volt} \times \SI{0,1}{\ampere} = \SI{12}{\watt},
  \qquad
  P_\text{média} = P_\mathrm{pico}\, (2\, t_p\, f) = \SI{12}{\watt} \times 0{,}03 = \SI{0,36}{\watt},
  \label{eq:ex-potencia}
\end{equation}
com $t_p = \SI{300}{\micro\second}$ e $f = \SI{50}{\hertz}$. A câmera térmica mostrou que o
aquecimento da revisão A não vinha dos MOSFETs, mas do regulador linear que alimentava os drivers da
ponte, dissipando \SI{1,4}{\watt}. Troquei-o por um regulador chaveado e pus vias térmicas sob cada
MOSFET.

\textbf{O limite de corrente.} Um comparador lê a tensão em $R_S$; acima de \SI{1,1}{\volt}
(\SI{110}{\milli\ampere}), ele aciona uma trava que desliga os drivers da ponte H, sem passar pelo
firmware. A trava só se desfaz por comando do firmware depois de a pessoa que opera apertar o reset,
e o botão de parada de emergência corta a habilitação do boost diretamente.

\textbf{O layout.} Quatro camadas: sinais e componentes em cima, terra contínuo na segunda, as
alimentações de \SI{12}{\volt} e de \SI{120}{\volt} em ilhas separadas na terceira, sinais embaixo
(Figura~\ref{fig:ex-layout}). Entre os dois lados, um rasgo de \SI{2,5}{\milli\meter} e
\SI{8}{\milli\meter} de distância de escoamento, com a tabela da IEC 60601-1 como referência. A malha de
corrente de cada canal ficou tão pequena quanto possível, com ligação Kelvin em $R_S$; o nó chaveado
da ponte ficou longe da entrada do encoder, que ganhou um filtro RC e um Schmitt trigger. A verificação
de regras (DRC) usou \SI{0,3}{\milli\meter} de isolação geral e \SI{1,5}{\milli\meter} na área de alta
tensão, e a verificação elétrica (ERC) terminou sem erro.

\begin{figure}[htbp]
  \centering
  % Planta simplificada do layout da revisão B (exemplo fictício): as áreas,
% a barreira de isolação e os estágios de saída. Não é o layout completo.
\begin{tikzpicture}[x=0.95mm, y=0.95mm, font=\scriptsize]
  % a placa, 100 x 80 mm
  \draw[line width=0.8pt, fill=nro-axon!8, rounded corners=1.5] (0,0) rectangle (100,80);
  \foreach \x/\y in {4/4,96/4,4/76,96/76} \draw[fill=white] (\x,\y) circle (1.6);
  % lado ligado ao USB
  \draw[densely dashed, nro-cinza, fill=nro-fundoclaro] (7,8) rectangle (33,72);
  \node[align=center, text=nro-cinza] at (20,67) {controle};
  \draw[fill=nro-cinza!45] (12,42) rectangle (26,56) node[midway, text=white] {MCU};
  \draw[fill=nro-cinza!25] (8,18) rectangle (16,24) node[midway, font=\tiny] {USB};
  \draw[fill=nro-cinza!25] (18,12) rectangle (30,22) node[midway, font=\tiny, align=center] {carga da\\bateria};
  % a barreira: o rasgo e os componentes que a cruzam
  \fill[white] (36,6) rectangle (38.5,74);
  \draw[nro-pulso-texto, line width=0.4pt] (36,6) rectangle (38.5,74);
  \draw[fill=nro-synapse!45] (32,28) rectangle (43,38) node[midway, font=\tiny, align=center] {DC-DC\\isolado};
  \draw[fill=nro-synapse!45] (32,48) rectangle (43,56) node[midway, font=\tiny, align=center] {isol.\\digitais};
  % lado do usuário (alta tensão)
  \draw[densely dashed, nro-pulso-texto, fill=nro-pulso!7] (46,8) rectangle (93,72);
  \node[text=nro-pulso-texto] at (69.5,67) {lado do usuário (120\,V)};
  \draw[fill=nro-tinta!55] (49,48) rectangle (63,60) node[midway, text=white, align=center, font=\tiny] {boost\\120\,V};
  \draw[fill=nro-tinta!30] (49,30) rectangle (61,40) node[midway, font=\tiny, align=center] {DAC e\\fontes};
  % os quatro estágios de saída, com as vias térmicas
  \foreach \i in {0,...,3} {
    \draw[fill=nro-tinta!75] (66+\i*6.5,14) rectangle (71+\i*6.5,22);
    \foreach \a in {0,1,2} \fill[nro-synapse!80!black] (66.8+\i*6.5+\a*1.6,12.2) circle (0.4);
  }
  \node[align=center, font=\tiny] at (79,27) {estágios de saída (Q1--Q4)\\e vias térmicas};
  % as saídas
  \foreach \i in {0,...,3} \draw[fill=nro-synapse!70] (95,14+\i*14) rectangle (100,22+\i*14);
  \node[anchor=west, align=left] at (101.5,40) {4 saídas\\(eletrodos)};
  % cotas
  \draw[{Stealth[length=1.4mm]}-{Stealth[length=1.4mm]}] (0,-4) -- (100,-4) node[midway, below] {100\,mm};
  \draw[{Stealth[length=1.4mm]}-{Stealth[length=1.4mm]}] (-4,0) -- (-4,80) node[midway, left] {80\,mm};
  \draw[nro-pulso-texto] (37.25,74) -- (37.25,86) node[above, align=center] {rasgo de 2,5\,mm: 8\,mm de\\distância de escoamento};
\end{tikzpicture}

  \caption{Planta simplificada do layout da revisão B: as áreas, a barreira de isolação e os estágios
    de saída. O layout completo está no repositório.}
  \fonte{Elaborado pela autora.}
  \label{fig:ex-layout}
\end{figure}
\end{desenvolvimento}

\begin{resultados}
As duas placas foram ensaiadas na bancada com o procedimento do plano de validação; a Atividade~4
descreve a montagem. A Tabela~\ref{tab:ex-revb} compara as revisões com os requisitos.

\begin{table}[htbp]
  \centering\small
  \caption{Revisão A e revisão B, contra os requisitos (média de 5 medições, placa 1).}
  \label{tab:ex-revb}
  \begin{tabular}{|l|c|c|c|}\hline
    \rowcolor{nro-fundo}\rotulo{Requisito} & \rotulo{Esperado} & \rotulo{Rev. A} & \rotulo{Rev. B} \\ \hline
    Amplitude em \SI{1}{\kilo\ohm}, \SI{100}{\milli\ampere} programados & $100 \pm 2$\,mA & 84\,mA & 100,9\,mA \\ \hline
    Amplitude em \SI{500}{\ohm}, \SI{100}{\milli\ampere} programados & $100 \pm 2$\,mA & 99\,mA & 100,4\,mA \\ \hline
    Temperatura do estágio de saída (60\,min, 50\,Hz) & $\le \SI{60}{\celsius}$ & \SI{71}{\celsius} & \SI{48}{\celsius} \\ \hline
    Ruído no sinal do encoder (pico a pico) & $\le \SI{50}{\milli\volt}$ & \SI{420}{\milli\volt} & \SI{35}{\milli\volt} \\ \hline
    Disparo do limite de corrente & 105 a 115\,mA & não havia & 108\,mA, em \SI{1,6}{\micro\second} \\ \hline
  \end{tabular}
  \fonte{Relatório de execução \texttt{NRO-PRO-003-901}.}
\end{table}

A primeira energização revelou dois defeitos, ambos meus: faltava o resistor de pull-down na
habilitação da trava, e as saídas podiam pulsar na partida; e o footprint do DC-DC isolado estava
espelhado. Os dois foram corrigidos à mão nas duas placas e estão entre os sete itens da errata que
virou a revisão C (\texttt{NBL-140}).
\end{resultados}

\begin{evidencias}
\evidencia{Repositório}{nebula-hw, tag rev-b}{O esquemático, o layout e a lista de materiais da revisão B, com os meus commits}{GitHub da equipe (acesso a pedido)}
\evidencia{Arquivo controlado}{NRO-PRO-014-901}{O design record da revisão B: requisitos, escolhas e alternativas}{SOMA › Arquivos}
\evidencia{Arquivo controlado}{NRO-PRO-003-901}{O relatório de execução dos ensaios de bancada da revisão B}{SOMA › Arquivos}
\evidencia{Atividade}{NBL-114; NBL-121; NBL-133}{Os cartões do esquemático, do layout e da montagem, atribuídos a mim e concluídos}{SOMA › Atividades › NEBULA}
\evidencia{Arquivo controlado}{NRO-PRO-004-901}{A USRS, de onde vêm os requisitos elétricos citados}{SOMA › Arquivos}
\end{evidencias}

\begin{aprendizados}
O curso me deu a fonte de corrente com amplificador operacional (Eletrônica Analógica), o
dimensionamento da dissipação (Eletrônica de Potência) e o transitório de primeira ordem (Circuitos
Elétricos II). O que ele não deu, e a placa deu: projetar para a segurança de uma pessoa ligada ao
circuito, ler uma norma de equipamento eletromédico, fazer layout pensando em ruído e em isolação, e
descobrir que a causa de um problema raramente é a primeira suspeita --- o calor não vinha de onde eu
esperava. Na próxima revisão, eu faria a revisão do footprint dos componentes novos contra o datasheet
antes de mandar fabricar, e não depois.
\end{aprendizados}
\end{atividade}
